% Generated draft: requires mathematical and source review.
\documentclass[11pt]{article}
\usepackage[margin=25mm]{geometry}
\usepackage{fontspec}
\setmainfont{Latin Modern Roman}
\usepackage{amsmath,amssymb,mathtools,unicode-math}
\setmathfont{Latin Modern Math}
\usepackage{enumitem,needspace,xurl,hyperref}
\setlength{\parindent}{0pt}
\setlength{\parskip}{5pt}
\setlength{\emergencystretch}{3em}
\newcommand{\N}{\mathbb{N}}
\newcommand{\Z}{\mathbb{Z}}
\newcommand{\Q}{\mathbb{Q}}
\newcommand{\R}{\mathbb{R}}
\newcommand{\C}{\mathbb{C}}
\newcommand{\F}{\mathbb{F}}
\newcommand{\A}{\mathbb{A}}
\newcommand{\PP}{\mathbb{P}}
\newcommand{\E}{\mathbb{E}}
\newcommand{\eps}{\varepsilon}
\newcommand{\dd}{\,\mathrm{d}}
\begin{document}
\section*{UnsolvedMath ID: 30004658}
\textbf{Source problem code:} catalogue entry
\subsection{Definitions and mathematical statement}
Write \(\mathbb T=\mathbb T_S(K)\) for the Hecke algebra named in the supplied record. Here \(S\) and \(K\) are source parameters whose meanings and permitted values are not specified. Let \(\operatorname{Max}(\mathbb T)\) denote its set of maximal ideals, and let \(\mathbb T_{\mathfrak m}\) denote localization at \(\mathfrak m\), under the intended commutative Hecke-algebra convention.

\textbf{Conjectural statement, with unresolved specifications.} For every
\[
\mathfrak m\in\operatorname{Max}(\mathbb T_S(K)),
\]
there exists a continuous semisimple residual Galois representation \(\bar\rho_{\mathfrak m}\) satisfying all three requirements named in the source:
\begin{enumerate}
\item its characteristic polynomials agree with the prescribed Hecke characteristic polynomials;
\item its complex conjugations are maximally odd;
\item it satisfies the prescribed Fontaine--Laffaille condition, described in the record as having Fontaine--Laffaille Hodge--Tate weights.
\end{enumerate}
If \(\mathfrak m\) is non-Eisenstein, the requested representation further admits a lift \(\rho_{\mathfrak m}\) with analogous properties, and the associated Taylor--Wiles deformation map is an isomorphism:
\[
R_S\longrightarrow\mathbb T_S(K)_{\mathfrak m}.
\]
The supplied statement places both the lifting assertion and the isomorphism assertion under the non-Eisenstein hypothesis.

\textbf{Limits of the formulation.} This is a schematic transcription of the requested assertions, not a fully specified conjecture. The record does not identify the underlying number field, Galois group, representation dimension, residual characteristic, coefficient fields or rings, or the topology defining continuity. It does not define the Hecke algebra, its cohomological realization, the Hecke characteristic polynomials, or the places where their comparison is required. It also does not define maximal oddness, the non-Eisenstein condition, the deformation problem represented by \(R_S\), or the displayed map. Hodge--Tate weights ordinarily concern characteristic-zero local representations; the residual Fontaine--Laffaille requirement therefore needs an explicit interpretation, such as an appropriate lift condition, which is not supplied here. No such interpretation is selected in this draft. The coefficient ring of the lift, its reduction to \(\bar\rho_{\mathfrak m}\), and the meaning of its analogous properties likewise remain unspecified.
\subsection{Short English statement}
The supplied conjecture asks whether every maximal ideal of the indicated Hecke algebra gives a residual Galois representation with prescribed Hecke, oddness, and Fontaine--Laffaille properties. For non-Eisenstein ideals it additionally asks for a lift and a Taylor--Wiles deformation-ring isomorphism. Essential arithmetic data and local conditions are missing, so this remains a schematic statement requiring mathematical review.
\subsection{Sources}
\begin{itemize}
\item Local UnsolvedMath registry, permanent ID \(30004658\), title \emph{Galois Representations from Hecke Maximal Ideals}; Oberwolfach Reports Open Problems collection; exact source problem code \texttt{catalogue entry}. The supplied record lists proposed year 2021 and names no individual proposer.
\item Supplied report citation: \emph{Arbeitsgemeinschaft: Derived Galois Deformation Rings and Cohomology of Arithmetic Groups} (2022), Oberwolfach Reports, pp.\ 1001--1046, DOI \texttt{10.4171/owr/2021/18}. \url{https://doi.org/10.4171/owr/2021/18}. The literature-source entry calls this \emph{Oberwolfach Report 18/2021}. These supplied year designations are retained without reconciliation.
\item Supplied literature source: Scholze, \emph{torsion in locally symmetric cohomology}. \url{https://arxiv.org/abs/1506.04144}.
\item Supplied literature source: Calegari--Geraghty, \emph{modularity beyond Taylor--Wiles}. \url{https://arxiv.org/abs/1209.6293}.
\item Literature assessment dated exactly \texttt{2026-08-22}: the record labels the problem partially solved and reports constructions of Hecke-eigensystem-valued residual Galois representations and many characteristic-zero or automorphic cases, while reporting that the full universal assertion with all listed local and lifting properties is not known in the stated generality. This is supplied dated curation derived from the OWR material, not an independently verified literature assessment; its time-sensitive status and cited claims require independent verification.
\item The supplied statement-status label is \texttt{corrected\_verified}. Its accompanying verification note states that only formula delimiters, notation, grammar, and typography were repaired, with mathematical content and scope preserved. This does not establish mathematical correctness or resolve the missing specifications above.
\end{itemize}
\end{document}
